\documentclass[11pt]{article}
\usepackage[margin=0.65in]{geometry}
\usepackage{graphicx}
\usepackage{amsmath}
\usepackage{newtxtext,newtxmath}
\usepackage{microtype}
\pagestyle{plain}
\begin{document}

\noindent\textbf{Figure 3: Fraction of Countries in R-Zone}

\smallskip
\noindent\footnotesize
The figure depicts the fraction of available countries in the R-Zone at a
given time:
\[
 \mathrm{Global\ R\mbox{-}Zone}_{t}
 = \frac{1}{N_t}\sum_{i\in S_t}\mathrm{R\mbox{-}Zone}_{it}.
\]
$N_t$ is the number of countries with the relevant debt and price-growth
measures in year $t$, and $S_t$ is the set of countries in the sample
at time $t$. Business combines business-debt growth with equity-price 
growth; household combines household-debt growth with house-price growth. 
R-Zone is the intersection of the highest debt-growth quintile and 
highest price-growth tercile. The figure uses the reconstructed 
historical panel for 1950--2012.

\begin{center}
\includegraphics[width=0.92\textwidth]{../_output/figure3_fraction_countries_rzone.jpg}
\end{center}

\end{document}
